% chapters/01_introduction.tex
\chapter{Introduction}\label{chap:introduction}

\section{Background}

The relationship between formal and informal sectors in developing economies has long been a central concern in development economics and public policy. India's dual economy presents a particularly compelling case study, with the informal sector employing over 80\% of the non-agricultural workforce \autocite{ILO2022}. Far from being a residual category of marginal workers awaiting absorption into modern employment, India's informal sector is deeply integrated into formal production networks through extensive subcontracting relationships.

The activities of registered mining, manufacturing, and non-manufacturing private or public enterprises under the Factories Act and/or the Companies Act that are involved in mining, manufacturing, and non-manufacturing activities are generally included in the organized (formal) segments of the economy in the Indian National Accounts Statistics (NAS). In 2019, the Ministry of Statistics and Programme Implementation (MoSPI) launched the Annual Survey of Unincorporated Sector Enterprises (ASUSE) to gather data on the unorganised sector, which is an essential but unseen component of India's economy. Its goal is to produce thorough, enterprise-level statistics on economic activity, employment, asset structure, and fundamental operational features of unincorporated businesses that are not included in formal registration regimes or the Annual Survey of Industries (ASI).

This persistence of informality challenges classical dual economy frameworks that predicted its inevitable decline. Economic development in dual economies has traditionally been theorized through a modernization lens, wherein the formal, modern sector inevitably absorbs or displaces the traditional, informal sector \autocite{Lewis1954}. This "displacement" narrative treats informality as a transitory friction—a vestige of pre-industrial production modes that modernization will eventually eliminate. Subsequent refinements incorporated migration costs \autocite{Harris1970}, credit constraints \autocite{Evans1989}, and regulatory barriers \autocite{Rauch1991}, but retained the core prediction: as institutions strengthen and capital deepens, informal enterprises either formalize or exit, unable to compete with formal sector productivity.

However, a fundamentally different interpretation emerges from structuralist political economy. Drawing on dependency theory and global value chain (GVC) analysis, a growing literature on formal-informal linkages rejects the dual economy perspectives and focuses on the nature of linkages and their effect on the developmental prospects of the informal economy \autocite{sethuraman1992urban, Harriss1990, Portes1989, Sassen1996, Meagher2010}. This "adverse incorporation" thesis holds that informality is not residual but \textit{functional}—systematically incorporated into capitalist production on highly unequal terms. Formal capital strategically deploys informal labor to extract surplus value while externalizing social reproduction costs. Far from competing with formal firms, informal producers serve as subordinated suppliers in vertically fragmented production chains, bearing volatility and regulatory risk while formal firms capture profits.

The Indian textile and apparel sector exemplifies this ambiguity. The sector accounts for approximately 7\% of manufacturing output and 12\% of export earnings, with employment exceeding 45 million workers \autocite{MinistryTextiles2023}. Production is organized through extensive vertical disintegration: large formal mills and exporters subcontract labor-intensive tasks—embroidery, stitching, dyeing, finishing—to networks of small unincorporated workshops and home-based workers. \textcite{Tewari1999} provides ethnographic evidence on the Ludhiana knitwear cluster, documenting how formal exporters coordinate production across hundreds of small job-workers without formal contracts. \textcite{Soundararajan2019} finds that informal wages rise when rural employment programs improve workers' outside options, suggesting binding monopsony power.



This ambiguity hinders effective policy formulation: if displacement dominates, policy should focus on enhancing informal sector productivity; if extraction dominates, regulation must address power imbalances in supply chains. To resolve the debate between the Displacement View (driven by technological efficiency) and the Extraction View (driven by unequal power), this dissertation maps these narratives to structural parameters—specifically the elasticity of substitution ($\rho$) and bargaining power ($\mu$).


\section{Research Gap}

Despite extensive descriptive research on formal-informal linkages, three critical gaps limit our understanding of the political economy of informality in India:

\textbf{First, the absence of structural identification.} Existing studies document correlations between formal growth and informal employment \autocite{Nataraj2011}, or between labor regulations and subcontracting intensity \autocite{Chaurey2015}, but cannot distinguish \textit{technological} complementarity from \textit{power-based} extraction. Observationally, both mechanisms predict similar outcomes: formal firms outsource, and informal wages remain low. Without a formal framework that separately identifies the elasticity of substitution ($\rho$) and bargaining power ($\mu$), policy recommendations remain ambiguous. Should interventions focus on enhancing informal productivity (if displacement dominates) or strengthening bargaining institutions (if extraction dominates)?

\textbf{Second, the treatment of labor institutions as exogenous constraints.} The labor economics literature on India \autocite{Besley2004, Ahsan2009} examines how statutory regulations affect formal employment but treats informal workers as a passive residual. This misses the strategic dimension: formal firms may \textit{actively respond} to enforcement by shifting production to the unregulated informal sector. The structural response—regulatory arbitrage—has first-order implications for understanding informality's persistence but remains empirically underexplored.

\textbf{Third, the lack of integrated formal-informal data.} Most studies analyze either formal sector data (ASI) or informal sector data (ASUSE/NSSO) in isolation, but not the \textit{vertical linkage} between them. The few exceptions \autocite{Singh2023} focus on internal contract labor rather than external outsourcing to separate enterprises. Quantifying pass-through elasticities—how productivity gains in formal firms translate (or fail to translate) into wages for informal suppliers—requires linked micro-data that existing studies lack.

This dissertation addresses these gaps by constructing the first state-industry panel linking formal sector productivity with informal sector wages and employment across a 15-year longitudinal horizon (2010--2024). By integrating NSS 67th, NSS 73rd, and ASUSE 21-24 rounds, I develop a CES-Nash bargaining framework that structurally distinguishes technology from power.

\section {Research Problem}


Existing literature lacks a quantifiable method to distinguish between these narratives. While studies show a link between formal and informal sectors, it remains unclear whether this link is driven by technological efficiency (displacement) or unequal bargaining power (extraction). Do formal firms outsource because informal suppliers are technologically complementary (as the displacement view suggests), or because fragmentation enables wage suppression (as the adverse incorporation view suggests)? When formal productivity rises, do informal wages increase through competitive pass-through, or do formal buyers capture the entire surplus through monopsony power? Does labor law enforcement protect informal workers by strengthening their bargaining position, or does it perversely drive production further underground?


\section{Research Questions}

The dissertation is organized around three core research questions:

\textbf{RQ1: Strategic Outsourcing}\\
\emph{Does formal sector productivity predict greater outsourcing to informal suppliers? If so, is this relationship stronger in states with stricter labor enforcement?}

\begin{itemize}
	\item \textbf{Hypothesis (Displacement):} Productive formal firms should \textit{reduce} outsourcing as they invest in modern internal technology (negative elasticity).
	\item \textbf{Hypothesis (Adverse Incorporation):} Productive formal firms should \textit{increase} outsourcing to minimize regulatory costs, especially in high-enforcement states (positive elasticity, stronger in high-enforcement environments).
\end{itemize}

\textbf{RQ2: Wage Pass-Through and Monopsony}\\
\emph{To what extent do productivity gains in formal firms pass through to wages of informal workers performing outsourced tasks?}

\begin{itemize}
	\item \textbf{Hypothesis (Competitive Markets):} Productivity gains should be bid into higher informal wages through competitive labor demand (positive pass-through elasticity).
	\item \textbf{Hypothesis (Monopsony/Extraction):} Formal buyers capture the entire surplus, with informal wages determined by subsistence constraints rather than productivity (zero pass-through elasticity).
\end{itemize}

\textbf{RQ3: Structural Embedding vs. Transition}\\
\emph{Do productive formal industries generate more informal enterprises (adverse incorporation) or absorb informal workers into formal employment (displacement)?}

\begin{itemize}
	\item \textbf{Hypothesis (Displacement):} Formal growth should reduce the count of informal enterprises (negative elasticity of informal firm density with respect to formal productivity).
	\item \textbf{Hypothesis (Adverse Incorporation):} Formal growth should generate informal satellites to perform subcontracted tasks (positive elasticity).
\end{itemize}


\section{Aim and Objectives}

\textbf{Aim:} To determine whether informality in India's textile manufacturing is a transitional phenomenon (displacement hypothesis) or a structurally embedded feature of capitalist production (adverse incorporation hypothesis), and to quantify the role of labor institutions in shaping surplus distribution between formal buyers and informal suppliers.

\textbf{Specific Objectives:}

\begin{enumerate}
	\item \textbf{Empirical Objective 1:} To estimate the elasticity of outsourcing with respect to formal sector productivity, and to test whether this relationship exhibits heterogeneity across states with varying labor enforcement intensity.
	
	\item \textbf{Empirical Objective 2:} To measure the pass-through of formal productivity gains to informal wages, distinguishing between competitive wage-setting (positive pass-through) and monopsonistic extraction (zero pass-through).
	
	\item \textbf{Empirical Objective 3:} To test whether productive formal industries generate informal "satellite" enterprises (adverse incorporation) or absorb informal workers into formal employment (displacement).
	
	\item \textbf{Theoretical Objective:} To develop a tractable structural model that formalizes the displacement-incorporation debate as competing specifications of technological substitutability ($\rho$) and bargaining power ($\mu$), and to derive testable predictions that distinguish between these mechanisms.


\end{enumerate}

\section{Significance of the Study}

This research makes four contributions to development economics, labor economics, and Indian economic policy:

\textbf{1. Empirical Contribution: Rigorous Quantitative Evidence on Surplus Distribution}

This is the first study to quantify pass-through elasticities between formal productivity and informal wages using a linked 15-year longitudinal panel ($N=152$). The finding that the elasticity is statistically zero ($p=0.636$) provides direct evidence for monopsonistic wage-setting in vertically fragmented value chains. Existing literature documents wage gaps descriptively \autocite{LaPorta2014} but does not test whether these gaps persist \textit{conditional on the value created}. By controlling for formal productivity, this study isolates the distributional mechanism from pure productivity differences.

\textbf{2. Theoretical Contribution: Formalizing the Displacement-Incorporation Debate}

The CES-Nash bargaining framework provides the first tractable structural model that nests both displacement and adverse incorporation as special cases. By mapping these theories to estimable parameters ($\rho$ for technology, $\mu$ for power), the framework transforms a qualitative theoretical debate into a quantitatively testable proposition. The comparative dynamics simulation further demonstrates how these parameters jointly determine equilibrium dynamics—outsourcing trajectories, wage convergence, welfare distribution—providing a blueprint for future empirical work with richer data.

\textbf{3. Policy Contribution: Unintended Consequences of Labor Enforcement}

The finding that formal firms in high-enforcement states exhibit 4.4 times greater outsourcing elasticity (7.0$\times$10$^{-6}$ vs 1.6$\times$10$^{-6}$) challenges conventional wisdom on labor law reform. Policymakers typically assume stricter enforcement unambiguously benefits workers by raising compliance. This research shows a perverse substitution effect: enforcement may drive production into the unregulated informal sector, where workers lack legal protections entirely. This has direct implications for India's ongoing debates on labor code consolidation and "ease of doing business" reforms.

\textbf{4. Methodological Contribution: Linking Micro-Surveys Across Sectors}

The data construction strategy—extracting state codes from ASI factory IDs, mapping ASUSE districts to states via NSS regions, and harmonizing variable definitions across survey years—provides a replicable template for future research. Public-use manufacturing surveys are typically analyzed in isolation; this study demonstrates how cross-sectoral linkages can be recovered despite geographic suppression in confidentiality-protected data.

\section{Scope of the Research}

\textbf{Geographic Scope:} The analysis covers 28 Indian states with valid enforcement data for the period 2010-24, aggregated to the state-industry level due to district code suppression in public microdata. The sample includes 152 pooled state-industry-year observations.

\textbf{Sectoral Scope:} The study focuses exclusively on textiles (NIC-13) and wearing apparel (NIC-14). This sectoral restriction is justified on three grounds:

\begin{enumerate}
	\item \textbf{Task divisibility:} Textile production involves discrete, separable tasks (embroidery, finishing, dyeing) that can be reallocated between factory floors and external workshops without fundamental technological reorganization. This makes the CES production function assumption empirically plausible.
	
	\item \textbf{Prevalence of subcontracting:} The textile sector exhibits the highest intensity of formal-informal linkages in Indian manufacturing. ASUSE data identifies 7,519 job-work enterprises in these sectors (compared to <500 in capital-intensive sectors like chemicals or metals), providing sufficient statistical power.
	
	\item \textbf{Policy relevance:} Textiles employ 45+ million workers, making findings directly relevant to labor policy debates. The sector also faces intense international competition (particularly from Bangladesh and Vietnam), making regulatory arbitrage economically salient.
\end{enumerate}

\textbf{Temporal Scope:} The primary analysis uses 2023-24 cross-sectional data (N=47). Robustness checks incorporate 2021-22 and 2022-23 to construct a three-year panel (N=141), though district-level variation is unavailable across years. The enforcement variable ($E_s$) is time-invariant, derived from 2023 Ministry of Labour annual reports.

\textbf{Analytical Scope:} The study focuses on three mechanisms:
\begin{enumerate}
	\item Strategic outsourcing (productivity $\rightarrow$ outsourcing intensity)
	\item Monopsonistic wage-setting (productivity $\rightarrow$ informal wages)
	\item Industrial density effects (productivity $\rightarrow$ count of informal enterprises)
\end{enumerate}

It does \textit{not} analyze:
\begin{itemize}
	\item Worker-level micro-behavior (requires matched employer-employee data unavailable in ASUSE)
	\item Firm entry/exit dynamics (cross-sectional data only)
	\item International trade linkages (requires customs data not in public domain)
	\item Non-wage labor outcomes (working conditions, hours, informality transitions)
\end{itemize}

\textbf{Methodological Scope:} The empirical strategy combines reduced-form OLS regressions with an empirically-grounded Monte Carlo simulation. The simulation uses parameters derived from the ASI-ASUSE data to visualize the dynamic implications of the cross-sectional findings.

\section{Organization of the Dissertation}

The remainder of the dissertation is organized as follows:

\textbf{Chapter 2: Literature Review} surveys three bodies of work: theoretical models of dualism and structural heterogeneity, empirical evidence on formal-informal subcontracting, and the labor institutions literature on enforcement and bargaining power.

\textbf{Chapter 3: Theoretical Framework} develops the CES-Nash bargaining model, derives equilibrium conditions, and formalizes the displacement-incorporation distinction as competing parameter restrictions.

\textbf{Chapter 4: Data and Methodology} describes the ASI-ASUSE data construction, presents the three empirical tests (outsourcing, wages, density), and outlines the comparative dynamics simulation procedure.

\textbf{Chapter 5: Results and Findings} reports the core empirical results: positive outsourcing elasticity ($9.07 \times 10^{-4}, p=0.043$), zero wage pass-through (-0.014, p=0.636), confirmed cost-shifting ($\beta = -7.23, p = 0.009$), heterogeneity by enforcement (4.4$\times$ larger effect in high-enforcement states), and simulation results showing qualitative consistency with adverse incorporation dynamics.

\textbf{Chapter 6: Discussion and Policy Implications} interprets the findings in relation to the displacement-incorporation debate, including the confirmed GVC price-squeeze mechanism and the implications of structural path dependence ($\beta=0.8503$). It discusses limitations such as structural breaks and geographic aggregation, and derives policy recommendations on enforcement design and supply chain regulation.

\textbf{Chapter 7: Conclusion} summarizes the key contributions and outlines directions for future research, including district-level analysis with restricted-access microdata and international comparative work.
